% ============================================================
%  cap01_introducao.tex
%  Capítulo 1: Por que as fontes de dados importam para
%              avaliação de políticas públicas?
% ============================================================

% ════════════════════════════════════════════════════════════
\chapter{Por que as fontes de dados importam para avaliação de políticas públicas?}
% ════════════════════════════════════════════════════════════

\begin{boxdestaque}{Neste capítulo}
A escolha da fonte de dados não é uma decisão técnica menor --- é um passo metodológico com consequências diretas sobre a validade das conclusões de qualquer avaliação. Este capítulo apresenta os princípios que orientam essa escolha, organizados em torno de três perguntas: (i) \textit{o que distingue dados bons de dados ruins para uma determinada pergunta avaliativa?}; (ii) \textit{como pensar a complementaridade entre diferentes tipos de fontes?}; e (iii) \textit{o que o avaliador precisa verificar antes de depositar confiança em qualquer base?} O capítulo também oferece uma visão panorâmica do ecossistema de dados disponíveis para avaliação de políticas públicas no Brasil.
\end{boxdestaque}

% ────────────────────────────────────────────────────────────
\section{A fonte de dados como escolha metodológica}
% ────────────────────────────────────────────────────────────

Toda avaliação de políticas públicas começa com uma pergunta. Mas entre a pergunta e a resposta existe uma etapa que frequentemente recebe menos atenção do que merece: a escolha dos dados. Essa escolha não é neutra. A fonte de dados selecionada determina o que pode ser medido, com que precisão, para quem e em que período. Quando a fonte é inadequada para a pergunta, nem a mais sofisticada estratégia de identificação causal produz resultados válidos.

Considere um exemplo simples. Um avaliador quer estimar o impacto de um programa de capacitação profissional sobre a empregabilidade dos participantes. Se usar a PNAD Contínua como fonte de dados, estará trabalhando com uma amostra representativa da população geral --- mas provavelmente insuficiente para identificar os próprios beneficiários do programa, cujo número pode ser pequeno. Se usar a RAIS, terá dados administrativos sobre o emprego formal de praticamente toda a força de trabalho formal brasileira --- mas perderá de vista os trabalhadores informais, que frequentemente são o público prioritário desses programas. Se conseguir acesso aos registros administrativos do programa, terá os beneficiários identificados individualmente --- mas precisará de uma fonte externa para construir o grupo de comparação.

Cada uma dessas fontes captura uma dimensão diferente da realidade, e a pergunta avaliativa determina qual delas --- ou qual combinação --- é mais adequada. Essa lógica de adequação entre pergunta e fonte é o fio condutor deste guia.

\begin{boxatencao}{Atenção --- Não existe fonte perfeita}
Toda fonte de dados tem limitações. O Censo Demográfico tem cobertura universal, mas é realizado apenas a cada dez anos. A PNAD Contínua é anual e tempestiva, mas tem tamanho de amostra insuficiente para análises no nível de municípios pequenos. Os registros administrativos têm cobertura ampla dos eventos que registram, mas cobrem apenas quem interage com o sistema --- excluindo populações que precisam de serviços mas não os acessam. Reconhecer as limitações das fontes que se usa não é fraqueza metodológica: é o primeiro passo para interpretar resultados com honestidade.
\end{boxatencao}

% ────────────────────────────────────────────────────────────
\section{Tipos de fontes de dados: uma primeira distinção}
% ────────────────────────────────────────────────────────────

Antes de discutir os critérios de seleção de fontes, é útil entender as grandes categorias em que as fontes disponíveis se organizam. O Capítulo~2 apresenta uma tipologia completa; aqui basta introduzir a distinção fundamental.

\subsection{Registros administrativos}

Registros administrativos são dados coletados como subproduto da operação de sistemas governamentais ou privados: cada empregado registrado na RAIS\footnote{Relação Anual de Informações Sociais, do Ministério do Trabalho e Emprego.}, cada internação no SIH\footnote{Sistema de Informações Hospitalares do SUS, gerido pelo DATASUS.}, cada aluno matriculado no Censo Escolar\footnote{Levantamento censitário anual conduzido pelo Inep junto às escolas de educação básica.}, cada beneficiário do Bolsa Família no CadÚnico\footnote{Cadastro Único para Programas Sociais do Governo Federal, instrumento de identificação e caracterização das famílias de baixa renda, gerido pelo MDS.}. Esses dados não foram coletados para fins de pesquisa, mas para gerir programas e sistemas.

As vantagens são claras: cobertura potencialmente universal da população atendida pelo sistema, baixo custo de coleta para o pesquisador (os dados já existem) e, frequentemente, possibilidade de acompanhamento longitudinal por meio de identificadores individuais (CPF, NIS, código de escola). A limitação central é que o dado registrado é aquele que o sistema precisa registrar --- não necessariamente o que o avaliador quer medir.

\subsection{Pesquisas amostrais}

Pesquisas amostrais --- como a PNAD Contínua, a POF e as pesquisas de saúde e vitimização --- coletam dados diretamente dos informantes (domicílios, pessoas, estabelecimentos) por meio de questionários estruturados aplicados a amostras representativas da população. Permitem medir variáveis que os registros administrativos não capturam: percepções, comportamentos, condições de saúde, consumo, experiências de vitimização.

A limitação central é o tamanho da amostra: pesquisas amostrais são projetadas para produzir estimativas representativas para determinados domínios geográficos e populacionais, e análises que extrapolam esses domínios produzem estimativas com erros padrão proibitivamente grandes.

\subsection{Censos}

Censos são levantamentos que cobrem toda a população de interesse, não uma amostra. O Censo Demográfico do IBGE, o Censo Escolar do Inep e o Censo Agropecuário são exemplos. Combinam a cobertura universal dos registros administrativos com a riqueza temática das pesquisas amostrais --- mas a custos de coleta muito elevados, o que os torna periodicamente irregulares.

% ────────────────────────────────────────────────────────────
\section{Critérios para avaliar a adequação de uma fonte}
% ────────────────────────────────────────────────────────────

A escolha da fonte procede em duas etapas. Primeiro, identifica-se o \textbf{tipo} de fonte potencialmente adequado à pergunta avaliativa --- registro administrativo, pesquisa amostral, censo, sistema setorial etc., conforme a tipologia do Capítulo~2. Em seguida, dentro do tipo escolhido, verifica-se se uma \textbf{fonte específica} é de fato adequada. Essa segunda etapa é organizada em torno de oito critérios.

\subsection{Cobertura geográfica}

A fonte cobre o território relevante para a pergunta? Uma pesquisa representativa para o Brasil como um todo pode não produzir estimativas confiáveis para municípios individuais ou para regiões específicas. A PNAD Contínua, por exemplo, produz estimativas trimestrais para o Brasil, grandes regiões, UFs e regiões metropolitanas --- mas não para municípios. Se a pergunta avaliativa requer análise municipal, é preciso buscar fontes com cobertura universal no nível do município (como o Censo Escolar ou a RAIS) ou construir estratégias de análise compatíveis com o que a fonte oferece.

\subsection{Periodicidade}

A periodicidade da fonte permite capturar o período relevante --- antes e depois da intervenção? Uma política implementada em 2018 que precisa ser avaliada com base em dados de 2016 e 2020 requer fontes que cubram esses dois momentos com metodologia comparável. Quebras metodológicas em séries históricas --- quando a fonte muda o questionário, a amostra ou os procedimentos de coleta entre as edições --- podem impossibilitar comparações ao longo do tempo.

\subsection{Unidade de análise}

A unidade de análise da fonte é compatível com a da pergunta? Se a pergunta é sobre alunos individuais, é preciso microdados de alunos (como o Censo Escolar), não dados agregados por escola ou município. Se a pergunta é sobre \textit{desempenho médio por escola}, tanto o SAEB quanto o Censo Escolar oferecem agregações adequadas --- mas agregar microdados de aluno do Censo Escolar preserva mais flexibilidade analítica do que trabalhar com médias pré-calculadas. Trabalhar com unidades de análise incorretas leva a falácias ecológicas e erros de inferência.

\subsection{Representatividade}

Para pesquisas amostrais: a amostra é representativa para o domínio de interesse? Verificar os domínios de representatividade declarados pela pesquisa antes de qualquer análise. Para registros administrativos: a cobertura do sistema é suficiente para os fins da análise, ou há exclusão sistemática de subgrupos relevantes?

\subsection{Granularidade temática}

A fonte contém as variáveis de que o analista precisa --- ou variáveis a partir das quais elas podem ser construídas? Um estudo sobre o impacto da escolaridade materna sobre saúde infantil precisa de uma fonte que contenha simultaneamente informações sobre a mãe e sobre a criança. Verificar os dicionários de variáveis antes de comprometer o desenho de avaliação com uma fonte específica.

\subsection{Identificadores para integração}

A fonte tem identificadores que permitem cruzar com outras bases de que o analista precisa? CPF, NIS, CNPJ, código de escola do Inep e código de estabelecimento de saúde (CNES) são os principais identificadores que viabilizam integrações no ecossistema de dados brasileiro. A disponibilidade desses identificadores --- e as restrições de acesso que os acompanham --- é uma dimensão crítica na fase de planejamento. Um exemplo concreto ilustra o valor desses identificadores: o código CNES identifica univocamente cada estabelecimento de saúde e está presente tanto no cadastro do CNES propriamente (com informações de infraestrutura e recursos humanos) quanto no SIH (com informações de internações e mortalidade hospitalar), permitindo, por exemplo, relacionar a oferta de leitos de UTI em cada hospital à mortalidade hospitalar observada --- uma análise impossível sem o identificador compartilhado.

\subsection{Qualidade do preenchimento}

Qual é a taxa de não-resposta para as variáveis de interesse? Há subregistro sistemático --- concentrado em determinadas regiões, grupos populacionais ou períodos? O SIM, por exemplo, tem subregistro de óbitos mais elevado no Norte e Nordeste; o SINAN tem subnotificação estrutural que varia muito por tipo de agravo. Ignorar esses padrões de subregistro diferencial leva a conclusões enganosas sobre desigualdades regionais e populacionais.

\subsection{Documentação e historicidade}

A fonte tem documentação metodológica que permita entender o que cada variável mede de fato? Há quebras metodológicas documentadas que afetam comparações ao longo do tempo? Fontes com documentação adequada reduzem o risco de erros de interpretação; fontes com histórico longo permitem análises de tendências e construção de contrafactuais mais robustos.

% CORRIGIDO: título encurtado para caber na caixa
\begin{boxdica}{Dica — Consulte sempre a documentação metodológica}
Todo grande produtor de dados brasileiro publica notas metodológicas, questionários e dicionários de variáveis. O IBGE publica documentação detalhada para a PNAD Contínua, o Censo Demográfico e todas as suas pesquisas. O Inep publica os cadernos de metodologia do Censo Escolar, do SAEB e do ENEM. O DATASUS disponibiliza os manuais de instrução de preenchimento do SIM, do SINASC e dos demais sistemas. Ler essa documentação antes de iniciar a análise é um investimento que poupa erros custosos.
\end{boxdica}

% ────────────────────────────────────────────────────────────
\section{O ecossistema de dados brasileiro: uma visão panorâmica}
% ────────────────────────────────────────────────────────────

O Brasil dispõe de um ecossistema de dados para avaliação de políticas públicas notavelmente rico, organizado em torno de três grandes pilares institucionais.

\textbf{O IBGE} é o produtor central do sistema estatístico nacional. Conduz o Censo Demográfico (decenal), a PNAD Contínua (trimestral/anual), a POF, o Censo Agropecuário e diversas outras pesquisas temáticas. Produz também as malhas territoriais e setores censitários que servem de base para análises espaciais. Para qualquer avaliação que precise de dados representativos da população brasileira em suas dimensões socioeconômicas, demográficas e de condições de vida, o ponto de partida é invariavelmente o IBGE.

\textbf{Os ministérios e sistemas setoriais} produzem os registros administrativos que documentam a operação das políticas públicas nas diferentes áreas. O Ministério da Saúde, via DATASUS, mantém o SIM, o SINASC, o SIH, o SINAN, o CNES e outros sistemas. O Ministério da Educação, via Inep, conduz o Censo Escolar, o SAEB e o ENEM. O Ministério do Trabalho gerencia a RAIS e o CAGED. O Ministério do Desenvolvimento e Assistência Social opera o CadÚnico e os dados do Bolsa Família. A Secretaria do Tesouro Nacional mantém o Siconfi/FINBRA. Cada um desses sistemas é descrito em detalhe nos capítulos temáticos deste guia.

\textbf{As agências reguladoras e outros produtores de dados especializados} cobrem setores específicos com dados que não estão disponíveis no sistema estatístico central: a ANEEL para energia elétrica, a ANATEL para telecomunicações, o INPE para dados de desmatamento, o TSE para dados eleitorais e de financiamento de campanha.

\begin{table}[H]
\centering
\caption{Visão panorâmica do ecossistema de dados para avaliação de políticas públicas}
\label{tab:ecossistema_geral}
\small
\begin{tabularx}{\textwidth}{@{} p{3.5cm} p{3cm} X @{}}
\toprule
\textbf{Pilar} & \textbf{Principais fontes} & \textbf{Uso central em avaliação} \\
\midrule
IBGE           & Censo Demográfico, PNAD-C, POF, Censo Agropecuário & Perfil populacional, condições de vida, representatividade nacional \\
Saúde (MS/DATASUS) & SIM, SINASC, SIH, SINAN, CNES & Mortalidade, nascimentos, internações, infraestrutura de saúde \\
Educação (Inep) & Censo Escolar, SAEB, ENEM & Matrículas, desempenho cognitivo, acesso ao ensino \\
Trabalho (MTE)  & RAIS, CAGED & Emprego formal, salários, rotatividade \\
Proteção social (MDS) & CadÚnico, Bolsa Família, BPC & Famílias vulneráveis, transferências de renda \\
Finanças públicas (STN) & Siconfi/FINBRA, Portal da Transparência & Gasto público, execução orçamentária \\
Agências reguladoras & ANEEL, ANATEL, ANS, ANTT & Infraestrutura setorial, cobertura, qualidade \\
Outros produtores de dados & TSE, INPE, SNIS, Ipea & Eleitoral, ambiental, saneamento, indicadores derivados \\
\bottomrule
\end{tabularx}
\fonte{FGV CLEAR, elaboração própria.}
\end{table}

% ────────────────────────────────────────────────────────────
\section{Complementaridade e integração entre fontes}
% ────────────────────────────────────────────────────────────

Uma das características mais importantes --- e mais subutilizadas --- do ecossistema de dados brasileiro é a possibilidade de \textbf{cruzar bases de dados} por meio de identificadores individuais. O CPF vincula pessoas físicas em qualquer base que o contenha. O NIS conecta beneficiários de programas sociais ao CadÚnico. O código de escola do Inep está presente no Censo Escolar e no SAEB. O CNES identifica estabelecimentos de saúde no SIH, no SIA e no SINAN.

Essas integrações permitem construir perfis multidimensionais que nenhuma fonte isolada poderia fornecer. Um exemplo clássico: cruzar o CadÚnico com o Censo Escolar pelo número de identificação do aluno permite analisar simultaneamente as condições de vida da família (renda, composição familiar, acesso a serviços) e o desempenho e a trajetória escolar das crianças. Cruzar a RAIS com dados de mortalidade do SIM pelo CPF permite estimar o impacto da formalização do emprego sobre expectativa de vida.

A capacidade de fazer essas integrações é, em muitos casos, o que torna avaliações causais rigorosas possíveis no Brasil. O Capítulo~13 discute as boas práticas e as obrigações legais que acompanham o uso de integrações de bases.

% ────────────────────────────────────────────────────────────
\section{Síntese do capítulo}
% ────────────────────────────────────────────────────────────

Este capítulo estabeleceu os fundamentos que norteiam todo o guia. Os pontos centrais são:

\begin{itemize}
  \item A escolha da fonte de dados é uma \textbf{decisão metodológica} com consequências diretas sobre a validade das conclusões: não existe fonte perfeita, mas existem fontes mais ou menos adequadas para cada pergunta avaliativa;

  \item As fontes se organizam em três grandes categorias --- \textbf{registros administrativos}, \textbf{pesquisas amostrais} e \textbf{censos} --- com perfis de vantagens e limitações complementares que o avaliador precisa conhecer;

  \item Oito critérios estruturam a avaliação de adequação de uma fonte: cobertura geográfica, periodicidade, unidade de análise, representatividade, granularidade temática, identificadores para integração, qualidade do preenchimento e documentação metodológica;

  \item O ecossistema brasileiro é organizado em torno do \textbf{IBGE}, dos \textbf{sistemas setoriais} dos ministérios e das \textbf{agências reguladoras e outros produtores de dados especializados} --- cada pilar com fontes distintas e complementares;

  \item A \textbf{integração entre bases} por identificadores individuais é um dos ativos mais poderosos do ecossistema brasileiro, viabilizando análises multidimensionais que nenhuma fonte isolada permitiria.
\end{itemize}

\medskip

O capítulo seguinte aprofunda a tipologia das fontes --- pesquisas amostrais, censos, registros administrativos, sistemas setoriais, dados fiscais, dados territoriais, fontes internacionais e APIs --- antes que os capítulos temáticos cubram fonte por fonte cada área de política pública.